\documentclass[t, 9pt]{beamer}

%% =====================================================================
%%  HOW TO CUSTOMISE THIS SLIDES TEMPLATE
%% =====================================================================
%
%  1. UNIVERSITY & DEPARTMENT: edit the "UNIVERSITY / PRESENTATION
%     CONFIGURATION" block in Preamble_slides.tex (point 11):
%     \UniversityName, \UniversityLogo, \DepartmentName.
%     Everything else in the preamble can stay untouched.
%
%  2. TITLE, AUTHOR, PLACE: edit these directly in
%     Titlepage_Options/Titlepage_slides.tex, inside the \title{},
%     \author{}, and \institute{} commands. Add extra lines (e.g. a
%     supervisor's address inside \institute{}) as needed.
%
%  3. CHAPTERS / AGENDA: edit \ChapterList in the config block --
%     it's a single comma-separated line, e.g.
%       \def\ChapterList{Intro,Theory,Data,Results,Conclusion,Appendix}
%     Add a name to insert a chapter, delete one to remove it. No other
%     file needs to change for this step.
%
%  4. STARTING A NEW CHAPTER IN THE SLIDES: use
%       \section{ChapterName}
%       \SetFootline{N}
%     where ChapterName matches an entry in \ChapterList, and N is that
%     entry's position (counting from 1). This highlights the current
%     phase in the progress bar at the bottom of every following slide.
%     IMPORTANT: if you add/remove a chapter in step 3, double-check that
%     all \SetFootline{N} numbers after that point still match the new
%     list order.
%
%  5. ADDING A SLIDE: use
%       \begin{frame}{Slide Title}
%           content here
%       \end{frame}
%     Add [c] after \begin{frame} to vertically center short content,
%     e.g. \begin{frame}[c]{Slide Title}.
%
%  6. TABLE OF CONTENTS (optional): Template_slides.tex has two commented-
%     out options right after the title page -- a single upfront TOC
%     slide, or a recurring outline slide shown before every section.
%     Uncomment ONLY ONE of the two blocks if you want either.
%
%  7. THEME COLOUR: change \definecolor{ThemeColor}{RGB}{...} in the
%     config block to re-brand the whole template (list markers,
%     progress bar, tcolorbox highlights all update automatically).
%
%  8. LOGO: swap the file at \UniversityLogo's path (config block) with
%     your own image, keeping the same filename or updating the path.
%
%  9. HIGHLIGHTED BOXES: wrap any key takeaway in
%       \begin{tcolorbox}[insightsbox]
%           Your text here.
%       \end{tcolorbox}
%     to get the pre-styled "Implication" box used throughout.
%
%  10. CITATIONS: cite sources on slides exactly like in the thesis, e.g.
%      \citep{miller2021} or \citet{miller2021}.
%
%  11. BIBLIOGRAPHY SLIDE: already included near the end of the file
%      (\bibliography{references}); no style needs to be set there.
%
%  Everything below THEME SETUP in Preamble_slides.tex is internal theme
%  logic (i.e. you shouldn't need to touch it for normal use).
%% =====================================================================



%%%%%%%%%%%%%%%%%%%%%%%%%%%%%%%%%%%%%%%%%%%%%%%%%%%%%%%%%%%%%%%%%%%%%%%%%%%%
% 1. FONTS: Times-style text and math font, matching the main document
%    exactly for visual consistency between thesis and slides.
%    Alternatives (swap this block):
%      Libertine:  \usepackage{libertine} + \usepackage[libertine]{newtxmath}
%      XCharter:   \usepackage{xcharter}
%%%%%%%%%%%%%%%%%%%%%%%%%%%%%%%%%%%%%%%%%%%%%%%%%%%%%%%%%%%%%%%%%%%%%%%%%%%%
\usepackage{newtxtext}

%%%%%%%%%%%%%%%%%%%%%%%%%%%%%%%%%%%%%%%%%%%%%%%%%%%%%%%%%%%%%%%%%%%%%%%%%%%%
% 2. MICROTYPOGRAPHY: subtle spacing tweaks for cleaner justified text.
%%%%%%%%%%%%%%%%%%%%%%%%%%%%%%%%%%%%%%%%%%%%%%%%%%%%%%%%%%%%%%%%%%%%%%%%%%%%
\usepackage{microtype}

%%%%%%%%%%%%%%%%%%%%%%%%%%%%%%%%%%%%%%%%%%%%%%%%%%%%%%%%%%%%%%%%%%%%%%%%%%%%
% 3. LANGUAGE & ENCODING: special characters + hyphenation rules.
%    Kept identical to the main document preamble so both compile
%    consistently if you copy content between them.
%%%%%%%%%%%%%%%%%%%%%%%%%%%%%%%%%%%%%%%%%%%%%%%%%%%%%%%%%%%%%%%%%%%%%%%%%%%%
\usepackage[utf8]{inputenc}
\usepackage[T1]{fontenc}
\usepackage[english]{babel}

%%%%%%%%%%%%%%%%%%%%%%%%%%%%%%%%%%%%%%%%%%%%%%%%%%%%%%%%%%%%%%%%%%%%%%%%%%%%
% 4. MATH: equations, symbols, theorem environments, bold vectors.
%    Identical setup to the main document so equations render exactly
%    the same way on slides as in the thesis.
%%%%%%%%%%%%%%%%%%%%%%%%%%%%%%%%%%%%%%%%%%%%%%%%%%%%%%%%%%%%%%%%%%%%%%%%%%%%
\usepackage{amsmath}
\usepackage{amssymb}
\usepackage{amsthm}
\usepackage{mathtools}
\usepackage{bm}
\usepackage{newtxmath} % load after amsmath; swap if using an alt font (see above)
\let\openbox\relax     % avoids a naming clash between amsthm and another package

%%%%%%%%%%%%%%%%%%%%%%%%%%%%%%%%%%%%%%%%%%%%%%%%%%%%%%%%%%%%%%%%%%%%%%%%%%%%
% 5. GRAPHICS, FIGURES & TABLES: same core packages as the main document,
%    minus the ones rarely needed on slides (subcaption, rotating, lscape).
%%%%%%%%%%%%%%%%%%%%%%%%%%%%%%%%%%%%%%%%%%%%%%%%%%%%%%%%%%%%%%%%%%%%%%%%%%%%
\usepackage{float}
\usepackage{caption}
\usepackage{booktabs}
\usepackage{array}
\usepackage{multirow}
\usepackage{adjustbox}

%%%%%%%%%%%%%%%%%%%%%%%%%%%%%%%%%%%%%%%%%%%%%%%%%%%%%%%%%%%%%%%%%%%%%%%%%%%%
% 6. TEXT DECORATION & LISTS
%%%%%%%%%%%%%%%%%%%%%%%%%%%%%%%%%%%%%%%%%%%%%%%%%%%%%%%%%%%%%%%%%%%%%%%%%%%%
\usepackage{enumitem}
\setlist[itemize,1]{label=\usebeamertemplate{itemize item}}
\setlist[itemize,2]{label=\usebeamertemplate{itemize subitem}}
\setlist[itemize,3]{label=\usebeamertemplate{itemize subsubitem}}

\usepackage{bbding}

%%%%%%%%%%%%%%%%%%%%%%%%%%%%%%%%%%%%%%%%%%%%%%%%%%%%%%%%%%%%%%%%%%%%%%%%%%%%
% 7. DIAGRAMS: TikZ for custom diagrams/flowcharts, identical setup to the
%    main document so figures can be copied between thesis and slides
%    without adjustment.
%%%%%%%%%%%%%%%%%%%%%%%%%%%%%%%%%%%%%%%%%%%%%%%%%%%%%%%%%%%%%%%%%%%%%%%%%%%%
\usepackage{tikz}
\usetikzlibrary{arrows.meta,positioning,fit,calc,decorations.pathreplacing}

%%%%%%%%%%%%%%%%%%%%%%%%%%%%%%%%%%%%%%%%%%%%%%%%%%%%%%%%%%%%%%%%%%%%%%%%%%%%
% 8. POSITIONING & COLORS: textpos allows placing elements at absolute
%    coordinates on a slide. xcolor must load before any \definecolor
%    call below, and before tcolorbox.
%%%%%%%%%%%%%%%%%%%%%%%%%%%%%%%%%%%%%%%%%%%%%%%%%%%%%%%%%%%%%%%%%%%%%%%%%%%%
\usepackage[absolute,overlay]{textpos}
\usepackage{xcolor}
\usepackage[most]{tcolorbox}

%%%%%%%%%%%%%%%%%%%%%%%%%%%%%%%%%%%%%%%%%%%%%%%%%%%%%%%%%%%%%%%%%%%%%%%%%%%%
% 9. BIBLIOGRAPHY: same author-year citation setup as the main document,
%    shared via the same file so \citep{}/\citet{} and the citation style
%    stay perfectly consistent between thesis and slides. Edit only in
%    Front Matter/bibliography_setup.tex.
%%%%%%%%%%%%%%%%%%%%%%%%%%%%%%%%%%%%%%%%%%%%%%%%%%%%%%%%%%%%%%%%%%%%%%%%%%%%
%%%%%%%%%%%%%%%%%%%%%%%%%%%%%%%%%%%%%%%%%%%%%%%%%%%%%%%%%%%%%%%%%%%%%%%%%%%%
% BIBLIOGRAPHY: author-year citations (natbib/plainnat = Harvard style,
%     common in economics).
%
%     HOW CITATIONS WORK (mechanism explained):
%     - All sources live in a separate plain-text file, e.g. references.bib,
%       NOT in this .tex file. Each entry has a unique key (e.g. smith2020)
%       plus fields like author, title, year, journal.
%     - \bibliography{references} at the END of your document (without the
%       .bib extension) tells LaTeX which file to pull from.
%     - \bibliographystyle{plainnat} below controls the FORMAT of both the
%       in-text citations and the final reference list. You never type the
%       full reference yourself; it's auto-generated from the .bib fields.
%     - \citep{smith2020}  ->  "(Smith, 2020)"   use when the name is NOT
%       part of your sentence grammar.
%     - \citet{smith2020}  ->  "Smith (2020)"    use when the name flows
%       into the sentence, e.g. "As \citet{smith2020} argues..."
%
%     ABBREVIATING A FREQUENTLY-CITED SOURCE:
%     If you cite the same paper repeatedly (e.g. a key theory paper),
%     define a shortcut once in the preamble:
%       \newcommand{\smithcite}{\citep{smith2020}}
%     Then just type \smithcite throughout instead of retyping the full
%     \citep{smith2020} each time. Useful if you later want to change how
%     that one source is cited everywhere at once.
%
%     COSMETIC APA-LIKE FALLBACK (stay on natbib/bibtex, no engine switch):
%     If a professor asks for "APA-ish" formatting but you don't want to
%     deal with biblatex+biber, just change the style below to:
%       \bibliographystyle{apalike}
%     Everything else (citep/citet, bibtex compilation) stays identical.
%
%     FULL APA 7 / CHICAGO AUTHOR-DATE (more accurate, bigger switch):
%     Requires replacing natbib entirely with biblatex + biber backend and
%     renaming citation commands (citep -> autocite/parencite,
%     citet -> textcite). Only worth it if the professor requires the
%     precise official style rather than something "close enough".
%%%%%%%%%%%%%%%%%%%%%%%%%%%%%%%%%%%%%%%%%%%%%%%%%%%%%%%%%%%%%%%%%%%%%%%%%%%%
\usepackage[round,authoryear]{natbib}
\bibliographystyle{plainnat}  % change to apalike for a quick APA-like look
\setlength{\bibhang}{1em}
\renewcommand{\bibfont}{\small}
\setlength{\bibsep}{1.2em}

%%%%%%%%%%%%%%%%%%%%%%%%%%%%%%%%%%%%%%%%%%%%%%%%%%%%%%%%%%%%%%%%%%%%%%%%%%%%
% 10. BEAMER BEHAVIOUR: hides the default navigation arrows/symbols at the
%     bottom of every slide.
%%%%%%%%%%%%%%%%%%%%%%%%%%%%%%%%%%%%%%%%%%%%%%%%%%%%%%%%%%%%%%%%%%%%%%%%%%%%
\beamertemplatenavigationsymbolsempty

%% =====================================================================
%% 11. UNIVERSITY / PRESENTATION CONFIGURATION
%%  --> Edit ONLY this block to adapt the template to a new university
%%      or department, and to re-brand the theme colour or agenda labels.
%%
%%  NOTE: The presentation title, author, and place are NOT set here.
%%  Edit them directly in Titlepage_Options/Titlepage_slides.tex, inside the
%%  \title{}, \author{}, and \institute{} commands. Add extra lines (e.g. a supervisor's address inside \institute{}) if required.
%% =====================================================================

% --- Institution (reused across all title page options and slide corners) ---
\newcommand{\UniversityName}{Your University}
\newcommand{\UniversityLogo}{images/university_logo.png}
\newcommand{\DepartmentName}{Department / Chair Name}

% --- Agenda / chapter labels used in the progress footline ---
\def\ChapterList{Chapter1,Chapter2,Chapter3,Chapter4,Chapter5,Appendix}

% --- Theme colour (change hex/RGB to re-brand) ---
\definecolor{ThemeColor}{RGB}{43,134,75}

%% =====================================================================
%%  THEME SETUP  (do not need to edit below this line for a new university)
%% =====================================================================

%%%%%%%%%%%%%%%%%%%%%%%%%%%%%%%%%%%%%%%%%%%%%%%%%%%%%%%%%%%%%%%%%%%%%%%%%%%%
% 12. THEME & COLOURS
%%%%%%%%%%%%%%%%%%%%%%%%%%%%%%%%%%%%%%%%%%%%%%%%%%%%%%%%%%%%%%%%%%%%%%%%%%%%
\usetheme{default}
\usecolortheme{Theme}

\setbeamertemplate{section in toc}[sections numbered]

%%%%%%%%%%%%%%%%%%%%%%%%%%%%%%%%%%%%%%%%%%%%%%%%%%%%%%%%%%%%%%%%%%%%%%%%%%%%
% 13. HIGHLIGHTED BOXES
%%%%%%%%%%%%%%%%%%%%%%%%%%%%%%%%%%%%%%%%%%%%%%%%%%%%%%%%%%%%%%%%%%%%%%%%%%%%
\tcbset{
  insightsbox/.style={
    enhanced,
    colback=ThemeColor!4,
    colframe=ThemeColor!85!black,
    boxrule=0.9pt,
    arc=2pt,
    left=6pt,right=6pt,top=6pt,bottom=6pt,
    fonttitle=\bfseries,
    coltitle=white,
    title={Implication},
  }
}

%%%%%%%%%%%%%%%%%%%%%%%%%%%%%%%%%%%%%%%%%%%%%%%%%%%%%%%%%%%%%%%%%%%%%%%%%%%%
% 14. FOOTLINE
%%%%%%%%%%%%%%%%%%%%%%%%%%%%%%%%%%%%%%%%%%%%%%%%%%%%%%%%%%%%%%%%%%%%%%%%%%%%
\usepackage{xstring}
\usepackage{etoolbox}
\newcounter{phasecount}
\newcounter{currentphase}
\setcounter{currentphase}{1}

\newcommand{\ProcessChapter}[1]{%
  \stepcounter{phasecount}%
  \ifnum\thephasecount=\thecurrentphase
    \textcolor{ThemeColor}{\textbf{#1}}%
  \else
    #1%
  \fi
  \ifnum\thephasecount=6\relax\else\ $\rightarrow$ \fi
}

%---- OPTIONAL: Exclude specific slides (e.g. title page, appendix,
%     bibliography) from the frame count shown in the footline ----%
% By default, the footline shows Beamer's built-in \insertframenumber and
% \inserttotalframenumber, which count every single frame in the document,
% including the title page, appendix, and bibliography slides.
%
% If you want an ACCURATE count that only includes your actual content
% slides, uncomment the three lines below AND manually add
% \stepcounter{realframe} inside every content frame that should count
% (i.e. every frame EXCEPT the title page, appendix, and bibliography).
%
% IMPORTANT: after activating this or changing which slides count, you
% need to compile TWICE for the total to display correctly. The first
% compile just records the count; the second one reads it back correctly.
% If the total looks wrong or blank, just recompile once more.
%
% \usepackage{totcount}
% \newtotcounter{realframe}
% \renewcommand{\insertframenumber}{\total{realframe}}
% \newcommand{\realtotal}{\total{realframe}}
% % Then replace \inserttotalframenumber with \realtotal in the footline
% % template below, and add \stepcounter{realframe} inside each counted frame.

\defbeamertemplate{footline}{progressfootline}{%
  \leavevmode\hbox{%
    \begin{beamercolorbox}[wd=.85\paperwidth,ht=2.5ex,dp=1.2ex,left]{}%
      \hspace{1em}\tiny
      \setcounter{phasecount}{0}%
      \edef\temp{\noexpand\forcsvlist{\noexpand\ProcessChapter}{\ChapterList}}%
      \temp
    \end{beamercolorbox}%
    \begin{beamercolorbox}[wd=.14\paperwidth,ht=2.5ex,dp=1.2ex,right]{}%
      \tiny\insertframenumber/\inserttotalframenumber\hspace{0.3em}%
    \end{beamercolorbox}%
  }%
}
\setbeamertemplate{footline}[progressfootline]

\newcommand{\SetFootline}[1]{%
  \setcounter{currentphase}{#1}%
}

\begin{document}

\title[Place]{Presentation Title}
\author[Author Short]{Author Name}
\institute[\UniversityName]{
     \\
    \DepartmentName, \\ \UniversityName
}
\vspace{2cm}

\titlegraphic{
     \begin{tikzpicture}[remember picture,overlay]
     \node[anchor=south east, yshift=8pt, xshift=0pt] at (current page.south east)
     {\includegraphics[width=1.5cm]{\UniversityLogo}};
     \end{tikzpicture}
 }

\addtobeamertemplate{frametitle}{}{%
\begin{tikzpicture}[remember picture,overlay]
     \node[anchor=south east, yshift=8pt, xshift=0pt] at (current page.south east)
     {\includegraphics[width=1.5cm]{\UniversityLogo}};
     \end{tikzpicture}
 }






\begin{frame}[plain]
  \titlepage
\end{frame}

%% =====================================================================
%%  TABLE OF CONTENTS OPTIONS (if wanted) 
%% =====================================================================

%---- OPTION A: Single TOC slide, shown once near the beginning ----%
% Simple and most common choice for a thesis defense: one outline slide
% right after the title page, not repeated anywhere else.

%\begin{frame}[plain]
%  \tableofcontents
%\end{frame}


%---- OPTION B: Recurring outline slide before every section ----%
% Shows a mini "where are we now" outline slide automatically before
% each \section{}, with the current section highlighted. Useful for
% longer presentations where the audience may lose track of structure,
% but adds one extra slide per chapter (e.g. 6 chapters = 6 extra slides).

%\AtBeginSection[]
%{
%  \begin{frame}[plain]
%    \frametitle{Outline}
%    \tableofcontents[currentsection]
%  \end{frame}
%}


%% =====================================================================
%%  This loads the tutorial for slides, comment out 
%%  when building your own and uncomment \section{Chapter1} till 
%%  \end{frame} to start your own. 
%%
%% =====================================================================

\section{Chapter1}
\SetFootline{1}

\begin{frame}{Scary Frame Environment}
Everything here works like in a regular DinA4 file. So you can use equations, tables, etc.\ like in the main template. There are just a few tweaks important for using the beamer template (that's the template for slides).

\vspace{0.3cm}
In general, slides require more case-by-case customisation by you, plus a lot of trial and error so it looks nice. 

\begin{itemize}
    \item Make use of \texttt{\textbackslash vspace\{\}} to create paragraph-like spacing.
    \item Use \texttt{\textbackslash small\{\}}, \texttt{\textbackslash footnotesize\{\}} etc. to scale any text.
\end{itemize}

\vspace{0.3cm}

\begin{enumerate}[label=\alph*)]
    \item If you use enumerate or itemize, think about specifying the label.
    \item Because in Beamer you get an error if you use enumerate and don't specify one (for options, see \texttt{introduction.tex}).
    \item For itemize lists I defined a standard here; adjust it if you want (this then overrides the standard locally).
\end{enumerate}

\end{frame}

\begin{frame}{Scary Frame Environment (2)}
    You can use the \texttt{columns} environment to get two columns side by side (adjusting each column's height with \texttt{\textbackslash vspace\{\}} again if one is shorter than the other):

\vspace{0.3cm}

\begin{columns}
  \begin{column}{0.48\textwidth}
    Each \texttt{column} takes a width (here \texttt{0.5\textbackslash textwidth}) and behaves like its own small block -- text, images, or tables can go inside.
  \end{column}
  \begin{column}{0.48\textwidth}
    The widths of all columns should add up to a bit less than 1 (e.g.\ two at 0.48 each) to leave a small gap between them.
  \end{column}
\end{columns}

\end{frame}



\begin{frame}{Basic Frame Types}
Every slide is its own frame. Unlike a DinA4 document, content does not flow automatically onto a next page -- if it overflows, it simply gets cut off, so each frame needs deliberate content limits.

\vspace{0.3cm}
\texttt{\textbackslash begin\{frame\}\{Title\} ... \textbackslash end\{frame\}}
\end{frame}

\begin{frame}[c]{Centralising Content: [c]}
Adding \texttt{[c]} right after \texttt{\textbackslash begin\{frame\}} vertically centers the content on the slide.

\vspace{0.3cm}
Useful for short content (one sentence, one equation, one research question) that would otherwise look sparse pinned to the top of an empty slide.
\end{frame}

\begin{frame}[plain]{Frame Without Footline: [plain]}
Adding \texttt{[plain]} removes the footline (progress bar, page number) for just this one slide, without affecting any other slide.

\vspace{0.3cm}
This is used for the title page and, if activated, the table of contents slide.
\end{frame}

\begin{frame}[fragile]{Frame With Special Content: [fragile]}
Adding \texttt{[fragile]} is required whenever a frame contains verbatim code (like the code snippets shown throughout this tutorial itself) or certain TikZ constructions.

\vspace{0.3cm}
If a slide throws an unexplained compilation error involving special characters, adding \texttt{[fragile]} is often the fix.
\end{frame}

\section{Chapter2}
\SetFootline{2}

\begin{frame}{How Chapters and the Progress Bar Work}

The progress bar you see at the bottom of this slide is customised. So a short explanation is in order.

\vspace{0.3cm}
Each \texttt{\textbackslash section\{ChapterName\}} should be followed immediately by \texttt{\textbackslash SetFootline\{N\}}, where \texttt{N} is that chapter's position in \texttt{\textbackslash ChapterList} (defined in the preamble).

\vspace{0.3cm}
This highlights the current phase in the progress bar shown at the bottom of every slide in that chapter. If you look at the bottom of this slide and compare it with the previous one, you will see that the highlighted part moved from Chapter1 to Chapter2 (you can rename them in the Preamble\_slides.tex document, of course).
\end{frame}

\begin{frame}{Adding or Removing a Chapter}
Edit \texttt{\textbackslash ChapterList} once, in the preamble's configuration block. All \texttt{\textbackslash SetFootline\{N\}} calls after the changed position must then be renumbered to match, since \texttt{N} refers to a chapter's position in the list, counting from 1.

\vspace{0.3cm}
(This is also explained in the relevant preamble section.)
\end{frame}

\begin{frame}{Cross-Referencing Slides}
\label{Slide 5}
This slide is labeled with \texttt{\textbackslash label\{Slide 10\}}. Later, writing "see Slide~\ref{Slide 5}" creates a clickable, auto-updating reference like in the main document.
\end{frame}

\begin{frame}[c]{Highlighted Takeaway Box}
\begin{tcolorbox}[insightsbox]
Key takeaways or implications go here.
\end{tcolorbox} 

\vspace{0.3cm}
This pre-styled box (colour, border, title) is already configured in the preamble -- just wrap any text you want to visually set apart.

\vspace{0.3cm}

specify the title locally if you don't want an Implication box using: \texttt{\textbackslash begin\{tcolorbox\}[insightsbox, title= Your Title]}.

\end{frame}

\begin{frame}{Shrinking a Wide Table}
Like in the main document \texttt{adjustbox} scales a table proportionally so it always fits the slide width without overflowing. For example: 

\vspace{0.3cm}

\begin{table}[t]
\centering
\begin{adjustbox}{max width=\textwidth}
\begin{tabular}{lcccccccc}
\toprule
 & (1) &  & (5) & (6) & (7) & (8) \\
 & 1985--1991 & [...] & 2008--2014 & \% & $\Delta$ 2014--1985 & \% \\
\midrule
Var mean log wages      & \textbf{61.7} &  & \textbf{94.6} & 100 & 32.9 & 100 \\
Var mean worker effects & 15.2 &  & 37.7 & 39.8 & 22.5 & \textbf{68.4} \\
Var mean plant effects  & 21.1 &  & 22.3 & 23.6 &  1.2 &  3.6 \\
Var mean $Xb$           &  0.8 &  &  0.6 &  0.6 & -0.2 & -0.7 \\
2 Cov(worker, plant)    & 19.3 &  & 39.4 & 41.7 & 20.1 & \textbf{61.0} \\
2 Cov(worker, $Xb$)     &  2.8 &  & -3.5 & -3.7 & -6.3 & -19.1 \\
2 Cov(plant, $Xb$)      &  2.4 &  & -1.9 & -2.0 & -4.3 & -13.1 \\
\bottomrule
\end{tabular}
\end{adjustbox}
\caption{Decomposition of across-city variation in average wages; see Table 3 in \citet{dauthMatchingCities2022}}
\label{tab:varDecomp}
\end{table}

\end{frame} %look into the file to see whats happening.


%% =====================================================================
%%  Your slides could start here. Comment out
%%  \section{Chapter1}
\SetFootline{1}

\begin{frame}{Scary Frame Environment}
Everything here works like in a regular DinA4 file. So you can use equations, tables, etc.\ like in the main template. There are just a few tweaks important for using the beamer template (that's the template for slides).

\vspace{0.3cm}
In general, slides require more case-by-case customisation by you, plus a lot of trial and error so it looks nice. 

\begin{itemize}
    \item Make use of \texttt{\textbackslash vspace\{\}} to create paragraph-like spacing.
    \item Use \texttt{\textbackslash small\{\}}, \texttt{\textbackslash footnotesize\{\}} etc. to scale any text.
\end{itemize}

\vspace{0.3cm}

\begin{enumerate}[label=\alph*)]
    \item If you use enumerate or itemize, think about specifying the label.
    \item Because in Beamer you get an error if you use enumerate and don't specify one (for options, see \texttt{introduction.tex}).
    \item For itemize lists I defined a standard here; adjust it if you want (this then overrides the standard locally).
\end{enumerate}

\end{frame}

\begin{frame}{Scary Frame Environment (2)}
    You can use the \texttt{columns} environment to get two columns side by side (adjusting each column's height with \texttt{\textbackslash vspace\{\}} again if one is shorter than the other):

\vspace{0.3cm}

\begin{columns}
  \begin{column}{0.48\textwidth}
    Each \texttt{column} takes a width (here \texttt{0.5\textbackslash textwidth}) and behaves like its own small block -- text, images, or tables can go inside.
  \end{column}
  \begin{column}{0.48\textwidth}
    The widths of all columns should add up to a bit less than 1 (e.g.\ two at 0.48 each) to leave a small gap between them.
  \end{column}
\end{columns}

\end{frame}



\begin{frame}{Basic Frame Types}
Every slide is its own frame. Unlike a DinA4 document, content does not flow automatically onto a next page -- if it overflows, it simply gets cut off, so each frame needs deliberate content limits.

\vspace{0.3cm}
\texttt{\textbackslash begin\{frame\}\{Title\} ... \textbackslash end\{frame\}}
\end{frame}

\begin{frame}[c]{Centralising Content: [c]}
Adding \texttt{[c]} right after \texttt{\textbackslash begin\{frame\}} vertically centers the content on the slide.

\vspace{0.3cm}
Useful for short content (one sentence, one equation, one research question) that would otherwise look sparse pinned to the top of an empty slide.
\end{frame}

\begin{frame}[plain]{Frame Without Footline: [plain]}
Adding \texttt{[plain]} removes the footline (progress bar, page number) for just this one slide, without affecting any other slide.

\vspace{0.3cm}
This is used for the title page and, if activated, the table of contents slide.
\end{frame}

\begin{frame}[fragile]{Frame With Special Content: [fragile]}
Adding \texttt{[fragile]} is required whenever a frame contains verbatim code (like the code snippets shown throughout this tutorial itself) or certain TikZ constructions.

\vspace{0.3cm}
If a slide throws an unexplained compilation error involving special characters, adding \texttt{[fragile]} is often the fix.
\end{frame}

\section{Chapter2}
\SetFootline{2}

\begin{frame}{How Chapters and the Progress Bar Work}

The progress bar you see at the bottom of this slide is customised. So a short explanation is in order.

\vspace{0.3cm}
Each \texttt{\textbackslash section\{ChapterName\}} should be followed immediately by \texttt{\textbackslash SetFootline\{N\}}, where \texttt{N} is that chapter's position in \texttt{\textbackslash ChapterList} (defined in the preamble).

\vspace{0.3cm}
This highlights the current phase in the progress bar shown at the bottom of every slide in that chapter. If you look at the bottom of this slide and compare it with the previous one, you will see that the highlighted part moved from Chapter1 to Chapter2 (you can rename them in the Preamble\_slides.tex document, of course).
\end{frame}

\begin{frame}{Adding or Removing a Chapter}
Edit \texttt{\textbackslash ChapterList} once, in the preamble's configuration block. All \texttt{\textbackslash SetFootline\{N\}} calls after the changed position must then be renumbered to match, since \texttt{N} refers to a chapter's position in the list, counting from 1.

\vspace{0.3cm}
(This is also explained in the relevant preamble section.)
\end{frame}

\begin{frame}{Cross-Referencing Slides}
\label{Slide 5}
This slide is labeled with \texttt{\textbackslash label\{Slide 10\}}. Later, writing "see Slide~\ref{Slide 5}" creates a clickable, auto-updating reference like in the main document.
\end{frame}

\begin{frame}[c]{Highlighted Takeaway Box}
\begin{tcolorbox}[insightsbox]
Key takeaways or implications go here.
\end{tcolorbox} 

\vspace{0.3cm}
This pre-styled box (colour, border, title) is already configured in the preamble -- just wrap any text you want to visually set apart.

\vspace{0.3cm}

specify the title locally if you don't want an Implication box using: \texttt{\textbackslash begin\{tcolorbox\}[insightsbox, title= Your Title]}.

\end{frame}

\begin{frame}{Shrinking a Wide Table}
Like in the main document \texttt{adjustbox} scales a table proportionally so it always fits the slide width without overflowing. For example: 

\vspace{0.3cm}

\begin{table}[t]
\centering
\begin{adjustbox}{max width=\textwidth}
\begin{tabular}{lcccccccc}
\toprule
 & (1) &  & (5) & (6) & (7) & (8) \\
 & 1985--1991 & [...] & 2008--2014 & \% & $\Delta$ 2014--1985 & \% \\
\midrule
Var mean log wages      & \textbf{61.7} &  & \textbf{94.6} & 100 & 32.9 & 100 \\
Var mean worker effects & 15.2 &  & 37.7 & 39.8 & 22.5 & \textbf{68.4} \\
Var mean plant effects  & 21.1 &  & 22.3 & 23.6 &  1.2 &  3.6 \\
Var mean $Xb$           &  0.8 &  &  0.6 &  0.6 & -0.2 & -0.7 \\
2 Cov(worker, plant)    & 19.3 &  & 39.4 & 41.7 & 20.1 & \textbf{61.0} \\
2 Cov(worker, $Xb$)     &  2.8 &  & -3.5 & -3.7 & -6.3 & -19.1 \\
2 Cov(plant, $Xb$)      &  2.4 &  & -1.9 & -2.0 & -4.3 & -13.1 \\
\bottomrule
\end{tabular}
\end{adjustbox}
\caption{Decomposition of across-city variation in average wages; see Table 3 in \citet{dauthMatchingCities2022}}
\label{tab:varDecomp}
\end{table}

\end{frame} and uncomment \section{Chapter1}
%%  till \end{frame} to start your own. 
%%
%% =====================================================================

%\section{Chapter1}
%\SetFootline{1}

%\begin{frame}{Slide Title}

%\end{frame}

%\section{Chapter2}
%\SetFootline{2}

%\begin{frame}{Slide Title}

%\end{frame}

\section{Chapter3}
\SetFootline{3}

\begin{frame}{Slide Title}

\end{frame}

\section{Chapter4}
\SetFootline{4}

\begin{frame}{Slide Title}

\end{frame}


\section{Chapter5}
\SetFootline{5}

\begin{frame}{Slide Title}

\end{frame}


\setbeamertemplate{footline}{}

\begin{frame}[c]{Bibliography}
    \tiny \bibliography{references}
\end{frame}

\section{Appendix}
\SetFootline{6}

\begin{frame}{Appendix: Slide Title}

\end{frame}

\end{document}


